\section{Experimental Setup}
\label{sec:setup}

\subsection{Data Splits}
\label{sec:splits}

We use a strictly temporal split to prevent any form of look-ahead bias:

\begin{table}[ht]
  \centering
  \caption{Temporal data splits.}
  \label{tab:splits}
  \small
  \begin{tabular}{@{}llc@{}}
    \toprule
    Split & Period & Fiscal period end date \\
    \midrule
    Train & 2012--2017 & $\leq$ 2017-12-31 \\
    Validation & 2018--2019 & $>$ 2017-12-31 and $\leq$ 2019-12-31 \\
    Test & 2020--2023 & $>$ 2019-12-31 and $\leq$ 2023-12-31 \\
    \bottomrule
  \end{tabular}
\end{table}

Splits are defined by the fiscal period end date (not the filing date or calendar year). This means a 10-K for fiscal year ending December 2017, filed in March 2018, is assigned to the training split. The test period deliberately spans the COVID-19 market disruption (2020) through the subsequent recovery, providing an out-of-distribution stress test.

All feature normalization statistics (medians, percentile bounds, quantile mappings) and coverage pruning decisions are computed on the training split only and applied identically to the validation and test splits.

Preliminary counts: ${\sim}$13{,}920 train / ${\sim}$5{,}657 val / ${\sim}$14{,}160 test observations. \todo{Verify exact split sizes from Colab notebook run.}

As a robustness check, we also define a walk-forward evaluation scheme using expanding training windows. Starting with a train cutoff of 2014-12-31, each successive fold advances by one year, with a 1-year validation window and a 2-year test window, up to a final test boundary of 2023-12-31. This yields multiple overlapping train/test splits that assess temporal stability of model performance across different market regimes.

\subsection{Evaluation Metrics}
\label{sec:metrics}

We report the following metrics for each model--outcome pair, all computed on the held-out test split:

\paragraph{Primary metric:} AUROC (Area Under the Receiver Operating Characteristic curve). This is the metric used for model selection, early stopping, and the go/no-go gate decision.

\paragraph{Secondary metrics:}
\begin{itemize}[leftmargin=*]
  \item \textbf{AUPRC} (Area Under the Precision-Recall Curve): more informative than AUROC under class imbalance, as it is sensitive to the positive class base rate.
  \item \textbf{Brier score}: mean squared error between predicted probabilities and binary outcomes, measuring both discrimination and calibration.
  \item \textbf{F1 at Youden threshold}: F1 score at the threshold maximizing Youden's $J$ statistic ($\text{TPR} - \text{FPR}$), providing a discrimination measure at a principled operating point.
  \item \textbf{ECE} (Expected Calibration Error): computed over 10 uniform-width bins, measuring the average gap between predicted probability and observed frequency.
\end{itemize}

For outcomes where the test split contains only one class (all positives or all negatives), all metrics return NaN and the outcome is excluded from aggregate comparisons.

In addition to aggregate metrics, we report AUROC and AUPRC stratified by Fama--French 12-industry sector. This sector-level breakdown surfaces whether model performance is concentrated in specific industries (e.g., finance) or generalises across the full company universe. Sectors with fewer than 30 observations or only one class in the test split are excluded from the stratified analysis.

\subsection{Go/No-Go Gate}
\label{sec:gate}

This study serves as a validation gate for subsequent work. The FT-Transformer passes the gate if it achieves AUROC at least 0.01 higher than XGBoost on three or more of the five distress outcomes (evaluated on the test set). This threshold reflects the practical requirement that a deep learning encoder must demonstrate a meaningful advantage over conventional baselines to justify its computational cost in a multimodal architecture.

\todo{The final benchmark (ATS-169) will include bootstrap confidence intervals (1,000 resamples) on AUROC differences to quantify statistical significance.}
